% ============================================================================
% The rule in the data (draft for the paper; Prompt 6.2, 2026-11-16).
% Numbers from scripts/10_ELSOC_rule.R (ELSOC 2016-2023, ponderador02, B = 500), the declared lambda
% sensitivity (05/06/06e with --tag=lambda_ext) and scripts/06i_q_format.R. Tables generated by
% 10_ELSOC_rule.R: output/10_ELSOC_rule/tab_rule_ladder.tex and tab_rule_forms.tex (SM table below);
% also available: tab_rule_E0.tex, tab_rule_threshold.tex, tab_rule_exposure_bins.tex,
% tab_rule_robustness.tex. Figure: plots/10_ELSOC_rule/rule_lambda_ladder.pdf.
% Full design note: docs/MAKING-IT-BOTTOM-UP.tex, section "PR #6: the rule in the data".
% Bibliography keys: oliver1993, manski1993, brock2001, shalizi2011, elsoc2023 (the last two added on
% 2026-11-16; check the citation the ELSOC release asks for).
% The 18-O pilot of Prompt 6.1 is NOT reported here (not robust; second paper): one sentence in the
% limitations at most, as in the last paragraph.
% ============================================================================
\subsection{The rule in the data}
\label{sec:rule-data}

The engine assumes a rule: an agent adopts when quality plus the weighted share of adopting
contacts reaches her aversion, $\Gamma + \lambda E_i \ge q_i$. In three of our four surveys
only $q_i$ is observed, so the form of the rule and the size of $\lambda$ are modelling choices.
ELSOC, the Chilean panel \citep{elsoc2023}, is the only one of the four surveys that measures
aversion, exposure and behaviour in the same people, and this is the role it plays in the
paper: it lets us estimate the rule. Respondents name the social movement they value most and
report how often, in the last 12 months, they took part in it ($y_i$: at least ``rarely''),
how often their family and their friends did ($E_i$: the mean of the two frequencies rescaled
to $[0, 1]$), and whether the movement's actions can bring social change (perceived efficacy,
our candidate for $\Gamma_i$). The willingness items that define $q_i$ are those of the
networks. We read the rule as a random-utility threshold model, $y_i = \mathbf 1[\Gamma_i +
\lambda E_i - q_i + \varepsilon_i \ge 0]$ with logistic noise of scale $\sigma$, the
specification of \citet{brock2001}. In the logit of $y_i$ on $\Gamma_i$, $E_i$ and $q_i$, the
rule is the restriction $\beta_\Gamma = -\beta_q$, and because $q_i$ is the engine's aversion,
$\hat\lambda = \beta_E/(-\beta_q)$ and $\hat\sigma = 1/(-\beta_q)$ are in engine units. All
estimates use the survey weights; intervals come from a bootstrap of respondents ($B = 500$).
The pooled sample has 5{,}192 person-waves from 3{,}331 respondents who named a movement (2016,
2018 and 2022, the waves with the willingness items).

\paragraph{Add, not multiply.} The additive rule fits far better than the product form used in
part of the diffusion literature: a pure product of quality and exposure loses 181 points of AIC
(138 with aversion measured a wave earlier) and has the worst cross-validated loss of every
specification with exposure. The product also predicts that nobody takes part without exposure,
whereas 14--15\% of respondents with $q_i \le 0.5$ whose family and friends never took part did.
Adding an interaction to the sum does not turn the rule into complements: the coefficient of
$\Gamma_i E_i$ is negative ($-3.8$ and $-3.0$, $p < 0.001$), so efficacy and exposure act as
substitutes. This agrees with a review of formal models of collective action: ``the multiplicative model rarely yields a significant improvement over a
linear model including the interest and efficacy terms'' \citep[p.~279]{oliver1993}.

\paragraph{How large is $\lambda$?} With the exposure respondents perceive in the same wave,
$\hat\lambda = 2.74$ [2.19, 3.66]: going from no exposure to full exposure offsets almost three
times the whole range of aversion. This is an upper bound. People who take part see their
friends take part, and with perceived exposure endogenous, contextual and correlated effects
cannot be separated \citep{manski1993}; latent homophily confounds contagion even in panels
\citep{shalizi2011}. We therefore climb a ladder of designs that remove part of the bias
(Table~\ref{tab:rule-ladder}). With lagged exposure and the lagged outcome in consecutive waves,
$\hat\lambda = 1.05$ [0.37, 2.64]. With a structural exposure built like the networks (for each
of the ego's 2017 confidants, the probability that a person with the confidant's attributes takes
part, averaged over confidants), which involves no report of the confidants' behaviour,
$\hat\lambda = 1.41$ [0.71, 2.50], and 0.77 [$-0.06$, 1.88] net of the ego's own attributes.
The data thus support $\lambda$ between about 0.7 and 1.5, with 2.5 as an upper bound, and
none of these estimates is causal. The range swept in the main analysis, 0.5--1.0, is the lower,
conservative part. At $\lambda = 1.25$ and $1.5$ (base exposure, random seeding, x2 networks)
the premium grows in all four surveys, to 9.5 (GSS 2004), 11.5 (ATP), 11.0 (GSS 2014) and 13.2
pp (ELSOC) at $\lambda = 1.5$. The region of the plane where the degree-preserving twin
overtakes the network is confined to $\lambda \le 1$ and disappears from $\lambda = 1.25$ in every
survey and under all five seeding strategies.

\paragraph{Quality is a common parameter.} Perceived efficacy carries no weight once aversion and
exposure are in the model ($\hat w_\Gamma = -0.08$ [$-0.36$, 0.28], against the 1 the rule
requires), and dropping it does not worsen the fit. Quality behaves as a property of the movement
and the moment, the same for everyone, which is how the engine treats $\Gamma$. Even across the
October 2019 outburst, with perceived exposure held fixed, $\Gamma$ moved by only $+0.06$
[$-0.06$, $+0.18$] in units of $q$; the rise in participation that year goes with exposure.

\paragraph{What the engine leaves out.} Two features of the estimated rule are absent from the
engine. The rule is noisy ($\hat\sigma \approx 0.5$ in units of $q$, including measurement
error), and exposure is concave: the first participating relative or friend moves participation
more than a linear index predicts, and full exposure moves it less. A flexible fit improves the
AIC by 195 points over the linear index, and the curvature lies in exposure, not in its
interaction with the distance to the threshold. A deterministic, linear rule therefore
overstates what dense exposure does. Both are natural extensions of the engine. The estimates
identify the form of the association between participation, aversion, efficacy and exposure in
the same people, not the causal effect of exposure. Only respondents who name a movement answer
these questions; weighting for that selection leaves $\hat\lambda$ practically unchanged (2.67).
Measuring aversion in the GSS format (whether the respondent has done each act, from the panel's
history, before asking about willingness) gives the same premium (7.65 against 6.95 pp,
$\lambda$-mean). Whether the networks also predict who joined the 2019 protests is a question
for future work: a first pilot was not robust to the noise of the rule.

\begin{table}[htbp]
\centering
\scriptsize
\setlength{\tabcolsep}{3pt}
\caption{The adoption rule estimated in ELSOC, 2016--2023. Panel A: $\hat\lambda$, the weight of
perceived efficacy and the noise of the rule, in units of aversion $q$, along the identification
ladder (weighted logits; 95\% intervals from a bootstrap of respondents, $B = 500$). Obs.:
person-waves; Resp.: respondents. S7 uses the 2017 confidants and the $q$ of 2016. S5: person
fixed effects in a linear probability model, reported as the ratio of the exposure coefficient
with and without them. Panel B: functional forms of the rule ($\Delta$AIC within sample; CV
log-loss: 10 folds by respondent; interaction: the coefficient of $\Gamma E$ or of $E q$).}
\label{tab:rule-ladder}
\textbf{A. How large is $\lambda$?}\\[2pt]
% generated by scripts/10_ELSOC_rule.R from output/10_ELSOC_rule/*.csv; do not edit by hand
\begin{tabular}{@{}llrrlll@{}}
\toprule
 & Specification & Obs. & Resp. & $\hat\lambda$ [95\% CI] & $\hat w_\Gamma$ [95\% CI] & $\hat\sigma$ \\
\midrule
\multicolumn{7}{@{}l}{\emph{Perceived exposure (family and friends in the same movement)}} \\
S1 & Pooled 2016, 2018, 2022; q of the same wave & 5\,192 & 3\,331 & 2.74 [2.19, 3.66] & $-$0.08 [$-$0.36, 0.28] & 0.48 \\
S2 & S1 + own attributes, movement, identification & 5\,094 & 3\,283 & 4.73 [3.32, 8.80] & $-$0.25 [$-$0.78, 0.31] & 0.80 \\
S3 & Lagged q (2016/18/22 $\to$ y of 2017/19/23) & 4\,794 & 3\,068 & 2.80 [2.13, 3.95] & 0.18 [$-$0.15, 0.56] & 0.53 \\
S3b & 2018 and 2019, q of 2018 & 4\,447 & 2\,969 & 2.45 [1.89, 3.22] & 0.05 [$-$0.29, 0.51] & 0.47 \\
S6 & Family only (friends in the model) & 5\,192 & 3\,331 & 1.20 [0.86, 1.69] & $-$0.08 [$-$0.34, 0.27] & 0.49 \\
S6 & Friends only (family in the model) & 5\,192 & 3\,331 & 1.56 [1.20, 2.19] & $-$0.08 [$-$0.34, 0.27] & 0.49 \\
\multicolumn{7}{@{}l}{\emph{Projection removed}} \\
S4 & Lagged E and lagged y, same movement & 1\,292 & 1\,016 & 1.05 [0.37, 2.64] & --- & 0.64 \\
S7 & Structural E of the 2017 confidants & 1\,242 & 1\,242 & 1.41 [0.71, 2.50] & --- & --- \\
S7 & \quad + own attributes & 1\,242 & 1\,242 & 0.77 [$-$0.06, 1.88] & --- & --- \\
S7 & \quad + own predicted probability & 1\,242 & 1\,242 & 0.65 [$-$0.17, 1.88] & --- & --- \\
\midrule
S5 & Person fixed effects (LPM): $\beta_E^{FE}/\beta_E^{pooled}$ & 10\,795 & 3\,956 & 0.80 ($\to\hat\lambda\approx$ 2.20) & --- & --- \\
\bottomrule
\end{tabular}
\\[8pt]
\textbf{B. Add or multiply?}\\[2pt]
% generated by scripts/10_ELSOC_rule.R from output/10_ELSOC_rule/*.csv; do not edit by hand
\begin{tabular}{@{}lrrrl@{}}
\toprule
Form & $\Delta$AIC & BIC & CV log-loss & Interaction \\
\midrule
\multicolumn{5}{@{}l}{\emph{Sample S1 (q of the same wave), n = 5\,192}} \\
\quad Additive (the rule) & 16.2 & 5306 & 0.4581 & --- \\
\quad Additive + $\Gamma E$ (complements) & 0.0 & 5297 & 0.4572 & $-$3.83 ($p < 0.001$) \\
\quad Pure product $\Gamma E$ & 197.0 & 5481 & 0.4735 & --- \\
\quad No exposure & 1564.7 & 6848 & 0.5926 & --- \\
\quad No efficacy & 14.6 & 5298 & 0.4577 & --- \\
\quad Additive + $E q$ & 8.6 & 5305 & 0.4579 & 1.80 \\
\multicolumn{5}{@{}l}{\emph{Sample S3 (lagged q), n = 4\,794}} \\
\quad Additive (the rule) & 8.6 & 5023 & 0.4805 & --- \\
\quad Additive + $\Gamma E$ (complements) & 0.0 & 5021 & 0.4804 & $-$2.96 ($p < 0.001$) \\
\quad Pure product $\Gamma E$ & 146.5 & 5154 & 0.4932 & --- \\
\quad No exposure & 1342.2 & 6350 & 0.6086 & --- \\
\quad No efficacy & 9.4 & 5017 & 0.4798 & --- \\
\quad Additive + $E q$ & 4.6 & 5025 & 0.4814 & 1.21 \\
\bottomrule
\end{tabular}

\end{table}
